% ============================================================
%  paper-export :: BRIEF 预设（简报 / 通报 / 会议纪要 / 情况说明）
%  配合: -V documentclass=ctexart
%
%  参照《党政机关公文格式》GB/T 9704-2012 的观感做的简化版：
%    无封面页、无目录、无页眉、无页码、不主动分页。
%    首页顶部是「文头」（标题 + 副标题 + 双分隔线），紧接着就是正文，
%    末尾右下角是落款（单位 / 编制人 / 日期）。
%
%  层级约定与 report 一致：H1 是文档名（被摘到文头），
%  H2 才是公文里的「一、」。
% ============================================================

% ---------- 1. 页面 ----------
% GB/T 9704 的版心：天头 37mm、地脚 35mm、左 28mm、右 26mm。
% 公文不装订成册，左边距不再额外留装订线。
\usepackage{geometry}
\geometry{a4paper, top=3.7cm, bottom=3.5cm, left=2.8cm, right=2.6cm}

\xeCJKsetup{CJKspace=true}

% 单字成行（「……确保审计可追 / 溯」这种末行只剩一个字）在中文排版里是硬伤，
% 仿宋三号一行才 28 个字左右，触发得比小四正文频繁得多。
% CheckSingle 让 xeCJK 主动把倒数第二行匀一个字下来，末行至少两字。
\xeCJKsetup{CheckSingle=true}

% ---------- 2. 字体 ----------
% macOS 的 STFangsong 位于系统字体集合中，xdvipdfmx 可能无法把它嵌入 PDF，
% 导致文件换到其他设备后正文空白。无封面版优先使用独立 OTF 的思源宋体，
% 缺失时回退到同样可嵌入的阿里巴巴普惠体；最后才使用系统仿宋。
\IfFontExistsTF{Noto Serif CJK SC}
  {\setCJKfamilyfont{zhfs}{Noto Serif CJK SC}
   \newfontfamily\pefangsong{Noto Serif CJK SC}
   \setmainfont{Noto Serif CJK SC}}
  {\IfFontExistsTF{AlibabaPuHuiTi_3_55_Regular}
    {\setCJKfamilyfont{zhfs}{AlibabaPuHuiTi_3_55_Regular}
     \newfontfamily\pefangsong{AlibabaPuHuiTi_3_55_Regular}
     \setmainfont{AlibabaPuHuiTi_3_55_Regular}}
    {\setCJKfamilyfont{zhfs}{STFangsong}
     \newfontfamily\pefangsong{STFangsong}
     \setmainfont{STFangsong}}}
\renewcommand{\fangsong}{\pefangsong\CJKfamily{zhfs}}

% 文头主标题：公文标准是「方正小标宋简体」—— 商业字体，装不了也不该依赖。
% Songti SC 自带 Black 字重，字形就是粗宋，是最接近小标宋的系统替代；
% 用黑体顶替则会丢掉宋体的横细竖粗，一眼就不像通报。
%
% macOS 自带的 Songti SC **有** Black 字重，但它藏在 Songti.ttc 这个字体集合
% 的非首 face 里，XeTeX/tectonic 取不到 —— 三种写法（family+BoldFont、
% 全名 Songti SC Black、PostScript 名 STSongti-SC-Black）实测全部落空，
% 静默回退到不含汉字的字体，标题整行排成豆腐块且编译不报错。
% （对比：\heiti 用的 Heiti SC Medium 能成，因为它是独立的 .ttf 而非集合成员。）
%
% 所以走思源宋体 Black（SIL OFL 免费，独立 OTF，字形正是粗宋）：
%   brew install --cask font-noto-serif-cjk-sc
% 没装则回退到 Songti SC 伪粗 —— 观感差一档但不会出豆腐块。
\IfFontExistsTF{Noto Serif CJK SC Black}
  {\setCJKfamilyfont{zhxbs}{Noto Serif CJK SC Black}
   \newfontfamily\pexiaobiaosong{Noto Serif CJK SC Black}}
  {\IfFontExistsTF{Source Han Serif SC Heavy}
    {\setCJKfamilyfont{zhxbs}{Source Han Serif SC Heavy}
     \newfontfamily\pexiaobiaosong{Source Han Serif SC Heavy}}
    {\setCJKfamilyfont{zhxbs}{Songti SC}[AutoFakeBold=2.5]
     \newfontfamily\pexiaobiaosong{Songti SC}[AutoFakeBold=2.5]}}
\newcommand{\xiaobiaosong}{\pexiaobiaosong\CJKfamily{zhxbs}}

% ---------- 3. 正文 ----------
% 三号（16pt）仿宋，行距固定 28pt —— 公文一页 22 行的来源。
% \linespread 是相对 1.2 倍基准的倍数：28 / (16 × 1.2) ≈ 1.46。
%
% 为什么要重定义 \normalsize：ctex 的 zihao 类选项只接受 {5, -4, 4}，
% 给它 3 会在 \ProcessKeysOptions 阶段直接报错、一页都出不来。
% export.sh 只能传四号打底，正文真正的三号靠这里补上 —— 而且必须走
% \normalsize，否则任何调用 \normalsize 的环境（列表、题注、脚注…）
% 都会把字号弹回四号。
\usepackage{indentfirst}
\renewcommand{\normalsize}{\zihao{3}}
\AtBeginDocument{%
  \setlength{\parindent}{2em}%
  \linespread{1.46}\fangsong\normalsize}
\setlength{\parskip}{0pt plus 1pt}

% ---------- 4. 标题层级 ----------
%   H1 → 三号黑体居中（正常情况下已被摘到文头，留着兜底）
%   H2 → 三号黑体      「一、」
%   H3 → 三号楷体      「（一）」
%   H4 → 三号仿宋加粗  「1.」
%   H5 → 三号仿宋加粗  「(1)」
% 公文的层级靠**字体**区分而非字号，四级标题与正文同为三号。
% 字号不拉开是刻意的 —— 通报正文短，用字号分级反而显得像技术手册。
%
% 标题一律左顶格（indent=0）。红头文件里标题确实和正文一样缩进 2 字，
% 但那套版式配的是满版心的红反线与文号栏，左边界另有锚点；这里只留了
% 双分隔线，标题再缩进就变成「线满宽、通篇无一行顶格」，左边白白空出
% 一条 1.2cm 的带子，版心看着整体右移了。缩进只留给正文段落首行。
\ctexset{
  section/format           = \heiti\zihao{3}\centering,
  section/beforeskip       = 12pt,
  section/afterskip        = 12pt,
  subsection/format        = \heiti\zihao{3},
  subsection/indent        = 0pt,
  subsection/beforeskip    = 8pt,
  subsection/afterskip     = 4pt,
  subsubsection/format     = \kaishu\zihao{3},
  subsubsection/indent     = 0pt,
  subsubsection/beforeskip = 6pt,
  subsubsection/afterskip  = 3pt,
  paragraph/format         = \fangsong\bfseries\zihao{3},
  paragraph/indent         = 0pt,
  paragraph/runin          = false,
  paragraph/beforeskip     = 4pt,
  paragraph/afterskip      = 2pt,
  subparagraph/format      = \fangsong\bfseries\zihao{3},
  subparagraph/indent      = 0pt,
  subparagraph/runin       = false,
  subparagraph/beforeskip  = 3pt,
  subparagraph/afterskip   = 2pt,
}

% 列表标记也顶格。preamble-common 把一级 itemize 的 labelindent 设成 2em
% 是为了和正文首行缩进对齐，在 brief 里会和标题的顶格打架 —— 同一份文档
% 里出现两条左边界，比哪条都不缩进更乱。
\setlist[itemize,1]{label=\textbullet,
                    labelindent=0em,labelwidth=0.9em,labelsep=0.5em,
                    leftmargin=1.4em}
\setlist[enumerate,1]{label=\arabic*.,
                    labelindent=0em,labelwidth=1.3em,labelsep=0.5em,
                    leftmargin=1.8em}

% ---------- 5. 分页 ----------
% preamble-common 给每级标题都 \pretocmd 了 \needspace{3.5~4.5\baselineskip}，
% 用意是「标题后至少留 3 行正文，别让标题孤零零挂在页面末行」。
%
% 那套阈值是按小四正文（行距 1.5 → \baselineskip ≈ 17.5pt）定的。brief 的
% \baselineskip 是 28pt，同样的「4.5 行」变成 126pt ≈ 4.4cm；简报又是标题密、
% 段落短的文体，几乎每个标题都会触发，结果是页面用到三分之一就硬断，
% 底下留一大片空白 —— 比标题挂在页底难看得多。
%
% \needspace 是 \pretocmd 上去的，没法「取消」，只能整体重定义：
% 收到 2.5\baselineskip（标题 + 1 行正文 + 余量，70pt），既挡住真正的孤标题，
% 又不会为了凑行数把半页内容推走。
%
% 2 行是不够的：标题自身的 beforeskip 也吃在这段空间里，恰好卡在边界时
% 标题会留在页底、正文被甩到下一页 —— 正是要防的情况。半行余量买断它。
\makeatletter
\AtBeginDocument{%
  \renewcommand{\needspace}[1]{\pe@needspace}}
\newcommand{\pe@needspace}{%
  \begingroup
    \setlength{\dimen@}{2.5\baselineskip}%
    \penalty-100\vskip\dimen@\penalty9999\vskip-\dimen@\penalty\z@
  \endgroup}
\makeatother

% 加粗段落（Markdown 里 `**第一步：准备演示环境**` 单独成段的写法）实际是
% 正文而非标题，LaTeX 不会为它做「后面至少留一行」的保护，于是它常常
% 孤零零躺在页底。给所有段落加一点 \clubpenalty 之外的保护成本太高，
% 这里只把段间断页的代价抬高，让 TeX 优先在别处断。
\interlinepenalty=200

% ---------- 6. 页眉页脚 ----------
% 一律为空。简报是「一口气读完」的文体，页码和页眉只会让它看起来像报告。
% \pagestyle{empty} 之外还得压掉 plain：article 的 \section* 会临时切到
% plain，否则首页底部会冒出一个孤零零的页码。
\makeatletter
\renewcommand{\ps@plain}{\let\@mkboth\@gobbletwo
  \let\@oddhead\@empty \let\@evenhead\@empty
  \let\@oddfoot\@empty \let\@evenfoot\@empty}
\makeatother
\pagestyle{empty}
\AtBeginDocument{\pagestyle{empty}\thispagestyle{empty}}

% ---------- 7. 图表题注 ----------
% 正文已是三号，题注用五号会显得断层，改用小四黑体。
\usepackage{caption}
\DeclareCaptionFont{pcap}{\heiti\zihao{-4}\color{pblack}}
\captionsetup{font=pcap,labelsep=quad,skip=6pt,justification=centering,
              singlelinecheck=true}
\captionsetup[figure]{position=bottom}
\captionsetup[table]{position=top}

% ---------- 8. 表格 ----------
% 正文三号偏大，表格照抄会撑爆版心。压到小四，行距回到 1.0。
% preamble-common 已把 \arraystretch 等设好，这里只调字号。
\AtBeginEnvironment{longtable}{\fangsong\zihao{-4}\linespread{1.0}\selectfont}
\AtBeginEnvironment{tabular}{\fangsong\zihao{-4}\linespread{1.0}\selectfont}
